\runninghead{Words That Do Not Pass}
\properlane{words-that-do-not-pass}{My Words Shall Not Pass: What God Has Said Is as Good as Done}

The Gospel ends with a parable and two sayings. ``From the fig tree learn a
parable: When the branch thereof is now tender and the leaves come forth, you
know that summer is nigh. So you also, when you shall see all these things,
know ye that it is nigh, even at the doors. Amen I say to you that this
generation shall not pass till all these things be done. Heaven and earth shall
pass: but my words shall not pass'' (Mt 24:32--35). The verse that follows,
which the Mass does not read, says that no one knows the day and the hour
(v.~36). Taken of the people then alive, the saying about this generation
would set a date for the end. The Fathers read here rest the certainty of the
promise on the one who makes it. The chants of the
Mass speak of deliverance in past tenses, \latin{avertísti}, ``thou hast
turned away,'' and \latin{Liberásti}, ``thou hast delivered,'' and the Epistle
thanks the Father who \latin{erípuit \dots\ et tránstulit}, ``hath delivered
\dots\ and hath translated'' us; yet the antiphon promises a return still to
come, and the Communion says \latin{crédite quia accipiétis}, ``believe that
you shall receive.'' St Augustine, on two of the psalms the Mass sings,
explains why a prophet speaks of the future as past. This interpretation
takes the past tenses and the Communion's future together: what God has said
is as certain as what he has done, and those who believe it pray as people
who will receive.

\subsection{The fig tree}
\label{sec:words-fig}

The Fathers do not read the fig tree alike, and this interpretation follows St
John Chrysostom and St Gregory the Great. Chrysostom says that the
parable shows the interval to be short, and more besides: it foretells ``a
spiritual summer, and a calm that should be on that day (after the present
tempest) for the righteous; but to the sinners the contrary, winter after
summer.'' Above all it is given ``to confirm His saying, as assuredly thus to
come to pass. For as this of the fig tree is of necessity, so that too,'' just
as the grain of wheat that must die proves the resurrection (Jn 12:24; 1~Cor
15:36; \work{Hom.\ in Mt.}~77.1).

Gregory, preaching on the same parable in Luke, where the fig tree and all the
trees ``shoot forth their fruit'' and the Lord says ``the kingdom of God is at
hand'' (Lk 21:29--31), reads it by contrast: \latin{sicut ex fructu arborum
vicina aestas cognoscitur, ita ex ruina mundi prope esse agnoscitur regnum
Dei}, ``as from the fruit of the trees the nearness of summer is known, so from
the ruin of the world the nearness of the kingdom of God is known.'' The
world's fruit is its fall: \latin{Ad hoc enim crescit, ut cadat}, ``For it
grows in order to fall.'' And the kingdom is rightly compared to summer,
\latin{quia tunc moeroris nostri nubila transeunt, et vitae dies aeterni solis
claritate fulgescunt}, ``because then the clouds of our grief pass away, and
the days of life shine with the brightness of the eternal sun''
(\work{Homiliae in Evangelia}~I.1.3).

St Jerome reads the parable as a warning against haste: when these things are
seen, \latin{nolite putare jam adesse consummationem mundi, sed quasi prooemia
et praecursores quosdam venire}, ``do not think that the consummation of the
world is already here, but that certain preludes and forerunners, as it were,
are coming'' (\work{In Matthaeum}~IV, on 24:32--33). St Augustine makes the fig
tree the human race, whose tender branch is the time when the sons of men
\latin{per fidem Christi ad spirituales fructus profecerint}, ``have advanced
through the faith of Christ to spiritual fruits'' (\work{Quaestiones
Evangeliorum}~I.39). St Hilary of Poitiers goes the other way and makes the
branch Antichrist, whose leafing is the green of sins in their triumph, and the
summer the heat of eternal fire (\work{In Matthaeum}~26.2).

\subsection{This generation}
\label{sec:words-generation}

Chrysostom will not take ``this generation'' of those then living. He says the
Lord speaks ``not of the generation then living, but of that of the believers,''
as the psalm says ``This is the generation of them that seek him'' (Ps 23:6),
and that ``over this generation shall nothing prevail, not famine, not
pestilence'' (77.1). Jerome gives two other answers: \latin{generatio} means
\latin{aut omne genus hominum \dots\ aut specialiter Judaeorum}, ``either the
whole human race, or the Jews in particular'' (on 24:34). Hilary speaks of
\latin{generationem nostram}, ``our generation,'' which cannot pass before all
these things have run their course, and he gives the purpose of the saying:
\latin{Atque ut fides certa esset futurorum, amen dicendo professione
veritatis, adjecit}, ``And so that faith in things to come might be certain,
he added, by saying Amen as a profession of truth,'' that this generation
would not pass (26.3).

\subsection{Heaven and earth, and the word}
\label{sec:words-word}

Gregory finds a paradox in the second saying: \latin{Nihil enim in
rerum corporalium natura coelo et terra durabilius, et nihil in rerum natura
tam velociter quam sermo transit}, ``For nothing in the nature of bodily things
is more lasting than heaven and earth, and nothing in nature passes so quickly
as speech.'' A word is not complete until it is spoken, and once it is spoken it
is gone. The Lord's meaning, he says, is that whatever seems lasting to us is
not lasting into eternity without change, \latin{quia sine mutabilitate
manentes sententias exprimit meus sermo qui transit}, ``because my speech,
which passes, utters sentences that remain without change''
(\work{Homiliae in Evangelia}~I.1.4).

Chrysostom says that heaven and earth are named because ``it were more easy for
these firm, fixed, and immoveable bodies to be blotted out, than for ought of
my words to fall to the ground,'' and to show the speaker to be ``maker of
all.'' Whoever doubts it should ``test His sayings, and when he hath found them
true \dots\ from what is past, let him believe also the things to come''
(77.1). Hilary gives the reason in the origin of each: heaven and earth,
\latin{profecta de nihilo}, ``come forth from nothing,'' \latin{habeant id in se
necessitatis ut non sint}, ``have in themselves the necessity of not being,''
while his words, \latin{ex aeternitate deducta}, ``drawn from eternity,''
\latin{id in se contineant virtutis ut maneant}, ``contain in themselves the
power to remain'' (26.3). Jerome adds that heaven and earth will pass
\latin{immutatione, non abolitione sui}, ``by change, not by their own
abolition'' (on 24:35), and St Bede, on Mark's form of the saying, that they
pass in the form they now have and subsist without end in their being (\work{In
Marci evangelium expositio}~IV, on 13:31). Hilary stresses that created things
can cease, Jerome and Bede that they will be changed; they differ in emphasis
and not in doctrine, and all of them set the word above the world.

Gregory begins the same sermon from the evidence of things already done. Many
of the evils foretold in the discourse he can see in his own time, and
\latin{cum multa praenuntiata jam completa sint, dubium non est quin sequantur
etiam pauca quae restant, quia sequentium rerum certitudo est praeteritarum
exhibitio}, ``since many things foretold are already fulfilled, there is no
doubt that the few that remain will follow, because the showing of past things
is the certainty of those that follow'' (I.1.1). Chrysostom made the same
argument at 77.1.

\subsection{The future told as past}
\label{sec:words-tenses}

St Augustine gives another ground, on the Gradual's own verse. Expounding
``thou hast saved us from them that afflict us,'' he writes: \latin{Et hoc
figura praeteriti de futuro dicitur: sed ideo tanquam praeteritum dicitur,
quia tam certum est quasi factum sit}, ``This too is said of the future in the
form of the past; but it is said as past because it is as certain as if it had
been done.'' The prophets, he explains, foretold the Passion in past tenses,
``They have dug my hands and feet'' (Ps 21:17). \latin{Nobis enim ea quae
praeterierunt, certa sunt; quae futura, incerta sunt}, ``For to us the things
that are past are certain, and the things to come uncertain.'' But \latin{Da
prophetam cui tam certum sit futurum quam tibi praeteritum}, ``Give me a
prophet to whom the future is as certain as the past is to you'': as what you
remember having happened cannot not have happened, so for him what he knows
will be cannot not be (\work{Enarrationes in Psalmos}~43.8). In the next verse
the psalm changes tense, \latin{laudábimur \dots\ confitébimur}, ``we shall
glory \dots\ we shall give praise,'' and Augustine notes it: \latin{Videte
quemadmodum miscet etiam verba futuri temporis}, ``See how he mixes in words
of the future tense too,'' so that the past tenses may be understood as
foretold of the future (43.9). On the Introit's verse he says the same: the
prophet \latin{utitur verbis quasi jam praeteriti temporis \dots\ quia apud
Deum et quod futurum est, jam factum est}, ``uses words as if of time already
past, because with God even what is to come is already done'' (84.3).

Chrysostom's argument and Augustine's are not the same. Chrysostom and
Gregory argue from what has already happened to what will happen; Augustine
argues from God, for whom the future is as fixed as the past.

The Epistle has both. Its petitions look forward: that the Colossians ``may be
filled,'' ``may walk,'' be strengthened ``in all patience and longsuffering.''
Its thanks look back to what the Father ``hath made,'' ``hath delivered'' and
``hath translated.'' Chrysostom notes that Paul, about to exhort them, ``makes
no mention of what by and by shall be laid up for them'' but ``mentions the
things which were already theirs, for these are the causes of the other.''
``For that which hath already come to pass gains belief'' (\work{Hom.\ in
Col.}~2). St Thomas Aquinas gives the translation into the kingdom two times:
\latin{Et hoc fit quando liberamur a peccato \dots\ Vel ad litteram, ut
consequeremur vitam aeternam}, ``And this happens when we are freed from sin
\dots\ or, according to the letter, that we might attain eternal life''; and he
takes \latin{longanimitas} as endurance against the delay of the reward,
citing Habacuc: ``if it make any delay, wait for it: for it shall surely
come'' (Hab 2:3; \work{Super Colossenses}, cap.~1, lect.~3).

\subsection{Promises, and those who ask}
\label{sec:words-introit}

The Introit puts the promise in God's mouth: \latin{Ego cógito cogitatiónes
pacis \dots\ invocábitis me, et ego exáudiam vos}, ``I think thoughts of peace
\dots\ you shall call upon me, and I will hear you.'' St Jerome, commenting on
the letter, says that God's thoughts are known to him alone, and that the
exiles' prayer was not needed for the promise to be kept: \latin{Certe absque
invocatione et oratione captivorum, poterat implere Dominus quod promiserat;
sed hortatur eos ad preces, ut quod promissum est, mereantur accipere}, ``The
Lord could certainly have fulfilled what he had promised without the calling
upon him and the prayer of the captives; but he urges them to prayer, so that
they may be worthy to receive what is promised,'' and he cites ``Ask, and it
shall be given you'' (Mt 7:7; \work{In Hieremiam}~V, on 29:10--14, PL~24
col.~893).

The Collect asks that the faithful ``may the more abundantly partake of Thy
mercies'' (Lasance). Its Latin asks for \latin{remédia}, ``remedies,'' that are
\latin{maióra}, ``greater,'' and so for more than has yet been given. The
Alleluia and the Offertory cry from the depths,
\latin{Dómine, exáudi oratiónem meam}, ``Lord, hear my prayer.'' St Hilary of
Poitiers says that God, who is everywhere, hears in
any case, \latin{sed orandus a nobis est}, ``but he is to be prayed to by
us''; \latin{Naturae suae est ut audiat; sed fidei nostrae est ut precemur
audiri}, ``It belongs to his nature to hear; it belongs to our faith to pray
that we be heard'' (\work{Tractatus super Psalmos}, on Ps 129). Aquinas, on
Paul's ``we cease not to pray,'' defines prayer as \latin{ascensus mentis in
Deum}, ``the ascent of the mind to God,'' and says that we come to it with
devotion \latin{non ut eum flectamus, sed ut nos erigamus in eum}, ``not to
bend him, but to raise ourselves to him'' (\work{Super Colossenses}, cap.~1,
lect.~3).

The Communion is the Lord's own promise to prayer: \latin{quidquid orántes
pétitis, crédite quia accipiétis, et fiet vobis}, ``whatsoever you ask when
you pray, believe that you shall receive, and it shall be done unto you.''
Its verb is future, \latin{accipiétis}, ``you shall receive.'' Augustine's
prophet holds what is to come as certain as what is past, and whoever believes
that he shall receive holds the thing asked for with that certainty before it
is given.

Bede, on the verse about forgiveness that follows the Communion's in Mark,
sets out who may ask what: \latin{Notanda distinctio deprecantium}, ``The
difference among those
who pray is to be noted.'' Whoever has perfect faith, \latin{quae per
dilectionem operatur}, ``which works through love,'' can move spiritual
mountains by prayer or even by command; those not yet perfect should ask that
their sins be forgiven, \latin{et absque dubio impetrabunt quod petunt, si
tamen in se peccantibus aliis primo ipsi dimittunt}, ``and without doubt they
will obtain what they ask, provided that they first forgive others who sin
against them'' (\work{In Marci evangelium expositio}~III, on 11:25).

Sicard of Cremona, a medieval bishop who wrote on the offices of the year,
reads the Gradual, the Offertory and the Communion as a sequence: those
freed from captivity give thanks in the Gradual, pray in the Offertory, and
\latin{suae petitionis, in communione: Amen dico vobis, recipiunt
certificationem}, ``receive in the
Communion, \emph{Amen dico vobis}, the assurance of their petition''
(\work{Mitrale}~VIII.25). In another chapter he says of the same chants that
\latin{Dominus in communione certificat de salute}, ``the Lord in the
Communion assures them of salvation'' (VIII.23). Rupert of Deutz had already
taken the Gradual for the thanks of those freed and the Offertory for their
supplication; the Communion as assurance is Sicard's own.

\subsection{Desire for what remains}
\label{sec:words-secret}

The Gospel says that heaven and earth will pass. The Secret asks that,
``being delivered from earthly desires, we may pass on to the enjoyments of
heaven'' (Lasance), and its Latin calls what is passed to \latin{cæléstia
desidéria}, heavenly ``desires.'' If heaven itself is to pass, desires that
are to last must be for the one who made it. Gregory says why the end of the
world is no grief to those who love God:
\latin{Qui ergo Deum diligunt ex mundi fine gaudere atque hilarescere
jubentur, quia videlicet eum quem amant mox inveniunt, dum
transit is quem non amaverunt}, ``Those who love God are bidden to rejoice and
be glad at the world's end, because they soon find him whom they love, while
that which they did not love passes away'' (\work{Homiliae in Evangelia}~I.1.3).

The Postcommunion has both tenses. It speaks of the sacrament \latin{quæ
súmpsimus}, ``which we have received,'' and asks that ``whatever is evil in
our hearts may be restored by its gift of healing'' (Lasance), with a verb
whose fulfilment is still to come, \latin{curétur}, ``may be cured.''

The deceivers of the Gospel's middle verses, who say that Christ is ``here''
or ``there,'' in the desert or in the inner rooms, offer a certainty of
another kind, one that can be pointed to. Against them the Lord sets the
lightning, whose coming, Chrysostom says, ``needs not a herald'' (76.3), and
his own warning: ``Behold I have told it to you, beforehand'' (Mt 24:25).

\subsection{A date, and an unconditioned promise}
\label{sec:words-difficulty}

Two difficulties stand in the way. The first is the saying about this
generation. Read of those alive when it was spoken, it would be a prediction
that failed, since they did not see the Son of man come on the clouds. None
of the Fathers read here takes it so, and their alternatives differ:
Chrysostom's generation of believers, Jerome's human race or Jewish people,
Hilary's ``our generation.'' Each of them leaves the day unknown, as
verse~36 says it is.

The second is the Communion. Read alone, it promises anything whatever to
anyone who believes. In Mark the promise is followed at once by the command to
forgive (Mk 11:25), and Bede, on that verse, puts the condition in the one who
prays: faith that works through love, and forgiveness of others. Three
sayings on other texts also look to the one who prays rather than to the thing
asked: Hilary's, on the psalm of the Alleluia and Offertory, that it belongs
to our faith to pray to be heard; Jerome's, on Jeremias's letter, that the
exiles are urged to pray so as to be worthy to receive what was promised; and
Aquinas's, on Paul's prayer for the Colossians, that we pray to raise
ourselves to God and not to bend him. None of the three comments on Mark.
Their agreement with Bede is this interpretation's reason for reading the
Communion as a lesson in how to pray, and not a guarantee of whatever object
is asked.

\begin{fourSenses}
\item[Literal.] Jesus tells his disciples that the signs will come as surely as
summer follows the fig tree's leaves, that this generation will not pass
before all is done, and that his words will outlast heaven and earth (Mt
24:32--35). The psalms of the Gradual and the Introit speak of deliverance as
done, and the Lord promises that what is asked in faith will be given (Mk
11:24).
\item[Allegorical.] The generation that does not pass is the generation of
believers, over which nothing prevails (Chrysostom); the word spoken in time
utters sentences that do not change (Gregory) and, drawn from eternity, has
the power to remain (Hilary); the prophet's past tense is the certainty with
God of what is to come (Augustine, 43.8, 84.3).
\item[Moral.] To believe the things to come from those already fulfilled
(Chrysostom, 77.1; Gregory, I.1.1); to pray as one who is promised a hearing,
since God hears by nature and faith asks (Hilary), and to forgive as one asks
to be forgiven (Bede); to wait for the reward that delays (Aquinas, on Hab
2:3).
\item[Anagogical.] The spiritual summer after the tempest for the just
(Chrysostom), when the clouds of grief pass and the days of life shine with
the eternal sun (Gregory); in this interpretation, the return the Introit
promises and the receiving the Communion promises are then complete.
\end{fourSenses}
